\chapter{Rotations, Spinors, and Color Symmetry}

\section{Rotations as a Group}

Rotations of space are symmetries, and they act on many kinds of objects. They
act on rigid bodies, on velocity vectors, on vector fields, and on quantum
states. An ordinary vector is the simplest place to begin because its
transformation can be pictured directly.

Rotations can also be combined. Rotate an object about one axis through one
angle, and then rotate it about a second axis through a second angle. The final
orientation is equivalent to a single rotation about some third axis through
some third angle. Thus the product of two rotations is another rotation.

This product is an operation on transformations, not an addition of the two
axis vectors. The first rotation changes the object on which the second acts.
If \(R_1\) is performed first and \(R_2\) second, the combined operation is
written \(R_2R_1\). It is itself some \(R_3\), even though finding the axis and
angle of \(R_3\) may be inconvenient. This closure under composition is the
first reason to regard rotations as a group rather than merely as a list of
geometrical motions.

It is useful first to describe one rotation geometrically. Choose a unit vector
\(\hat{\boldsymbol n}\) along the rotation axis and an angle \(\theta\). We
adopt the right-hand convention. A left-handed rotation through \(\theta\) can
equally be written as a right-handed rotation through \(-\theta\), or as a
right-handed rotation about \(-\hat{\boldsymbol n}\).

\begin{figure}[htbp]
\centering
\includegraphics[width=0.58\textwidth]{lecture_03_figure_02.png}
\caption{An axis \(\hat{\boldsymbol n}\) and an angle \(\theta\) specify a spatial rotation.}
\lectureframe{00:02:34}
\end{figure}

A direction on the unit sphere requires two parameters, such as longitude and
latitude. The angle about that direction supplies a third. Rotations therefore
form a three-parameter family. We shall not yet worry about the global
identifications caused by angles differing by \(2\pi\); the local parameter
count is what matters here.

The axis-angle description has some redundancy. Reversing the axis while
reversing the angle gives the same operation, and sufficiently large angles
can be reduced modulo a full turn. None of that changes the fact that near the
identity a rotation has three continuously variable parameters. That local
count will later be matched against the number of independent parameters in a
matrix representation.

The group properties can now be checked in turn:
\begin{enumerate}
\item \emph{Closure.} Performing two rotations in succession produces another
rotation.
\item \emph{Identity.} Doing nothing is the identity rotation.
\item \emph{Inverse.} A rotation through \(\theta\) about
\(\hat{\boldsymbol n}\) is undone by a rotation through \(-\theta\) about the
same axis.
\item \emph{Associativity.} For three successive rotations,
\((R_1R_2)R_3=R_1(R_2R_3)\).
\end{enumerate}

Rotations are not commutative. A rotation about the \(x\)-axis followed by one
about the \(y\)-axis generally gives a different orientation from the same two
operations in the reverse order:
\begin{equation}
R_x(\alpha)R_y(\beta)\ne R_y(\beta)R_x(\alpha).
\end{equation}
The difference becomes second order when both angles are very small, which is
why infinitesimal rotations can look almost commutative. The full group is
nevertheless non-abelian.

The noncommutativity can be checked without any equations. Take a book or a
box, mark one face as the front, and turn it by ninety degrees about the
\(x\)-axis and then by ninety degrees about the \(y\)-axis. Return to the
original orientation and reverse the order. The marked face ends in a
different direction. For small angles the discrepancy is harder to see, but
it has not vanished; it begins at the product of the two small angles.

\begin{classroomqa}
\audiencequestion{Are these rotations sequential rather than simultaneous?\lecturetimestamp{00:04:45}}
\lecturerresponse{Yes. Group multiplication means one operation followed by another. The order of those successive rotations is precisely what can fail to commute.}
\end{classroomqa}

This is our first non-abelian group. The eventual application is not the color
of an ordinary object but the red, green, and blue color states of a quark.
That is a three-state internal system, and its interesting transformations mix
the three colors. Ordinary rotations provide a familiar model of how a
continuous group acts through matrices.

A group by itself is abstract. To use it in physics we must say what objects it
acts on. A emph{representation} assigns an operator or matrix to each group
element so that multiplication of group elements is reproduced by
multiplication of the assigned matrices. Different physical objects can carry
different representations of the same group. Vectors, spin-one states, and
spin-one-half states will provide three examples.

\section{The Three-Dimensional Vector Representation}

Let an ordinary vector have components
\begin{equation}
\boldsymbol V=(V_x,V_y,V_z),\qquad V_i,\quad i=1,2,3.
\end{equation}
One may rotate the physical vector while holding the axes fixed, or rotate the
coordinate axes while holding the vector fixed. The conventions differ, but
in either description the components are mixed by a \(3\times3\) matrix.
For every axis-angle pair \((\hat{\boldsymbol n},\theta)\), write
\begin{equation}
V_i'=R_{ij}(\hat{\boldsymbol n},\theta)V_j,
\label{eq:l3-vector-rotation}
\end{equation}
where a repeated index is summed.

The repeated index \(j\) is a dummy label. Equation~\eqref{eq:l3-vector-rotation}
means explicitly
\begin{equation}
V_i'=R_{i1}V_1+R_{i2}V_2+R_{i3}V_3
\end{equation}
for each output component \(i\). The free index \(i\) labels which component
is being computed; the repeated \(j\) is summed and may be renamed without
changing the equation. Keeping those two roles separate prevents most index
errors in the length-preservation argument.

\begin{figure}[htbp]
\centering
\includegraphics[width=0.68\textwidth]{lecture_03_figure_03.png}
\caption{A rotation matrix maps the original vector components to the rotated components.}
\lectureframe{00:07:41}
\end{figure}

A true rotation preserves length. For the moment both \(V_i\) and \(R_{ij}\)
are real, so
\begin{equation}
V_i'V_i'=V_iV_i
\end{equation}
for every vector. To expose the matrix condition, write the same
transformation with different dummy indices:
\begin{equation}
V_i'=R_{ij}V_j,\qquad V_i'=R_{ik}V_k.
\end{equation}
Multiplying and summing over \(i\) gives
\begin{align}
V_i'V_i'
 &=R_{ij}V_jR_{ik}V_k\\
 &=R_{ij}R_{ik}V_jV_k.
\end{align}
This equals \(V_jV_j\) for every choice of \(V\) only if
\begin{equation}
R_{ij}R_{ik}=\delta_{jk}.
\label{eq:l3-orthogonality-index}
\end{equation}
In matrix notation,
\begin{equation}
R^TR=\boldsymbol 1,
\label{eq:l3-orthogonal}
\end{equation}
and hence
\begin{equation}
R^{-1}=R^T.
\end{equation}
Transposition simply interchanges rows and columns.

Why must the coefficient be \(\delta_{jk}\)? The Kronecker delta sets \(j=k\)
under the sum:
\begin{equation}
\delta_{jk}V_jV_k=V_1^2+V_2^2+V_3^2.
\end{equation}
If \(R_{ij}R_{ik}\) differed from \(\delta_{jk}\), some direction or some
cross term \(V_jV_k\) would change the length. Since the equality is required
for every possible vector, not merely for one chosen vector, the matrix
coefficient itself must equal the identity.

Equation~\eqref{eq:l3-orthogonal} is the defining orthogonality condition for
this real matrix representation. Multiplying the matrices reproduces the
multiplication table of spatial rotations. This is not the only
representation of the rotation group; it is merely the one most easily
visualized.

The inverse relation has an immediate geometrical meaning. A rotation can be
undone without solving a new matrix equation: exchange rows and columns. If
\(R_1\) and \(R_2\) represent two rotations, then \(R_2R_1\) represents their
successive action and is again orthogonal. Thus the matrix operations preserve
the same closure, identity, inverse, and associative structure established
geometrically.

\section{Spin One and the Vector Representation}

The three-dimensional representation also acts on quantum states. A spin-one
particle has three possible spin projections,
\begin{equation}
m=+1,\quad 0,\quad -1.
\end{equation}
Its spin state is therefore a three-component column of probability
amplitudes. Rotating the particle, or changing the orientation of a magnetic
field used to define the spin axis, mixes those three amplitudes by a
three-dimensional representation of the rotation group.

The number of spin states follows the familiar rule \(2s+1\). Spin zero has
one state, spin one-half has two, spin one has three, and spin three-halves has
four. The dimension of the state vector therefore changes with the spin. The
same physical rotation must be represented by matrices of different sizes
when it acts on these different state spaces.

For a spin-one particle, choose the basis
\begin{equation}
\lvert 1,+1\rangle,qquad
\lvert 1,0\rangle,qquad
\lvert 1,-1\rangle.
\end{equation}
A state is
\begin{equation}
\lvert\psi\rangle
=c_{+}\lvert1,+1\rangle+c_0\lvert1,0\rangle+c_{-}\lvert1,-1\rangle.
\end{equation}
Rotating the apparatus does not create a fourth spin state; it mixes the three
coefficients \(c_+,c_0,c_-\) among themselves. The corresponding rotation
operator must therefore be a \(3\times3\) matrix.

For spin one, this representation transforms in the same way as an ordinary
spatial vector. This is the precise meaning of calling a spin-one particle a
\emph{vector particle}: its three spin amplitudes and the three components of
a vector obey the same rotation law. The quantum state need not be imagined
as a literal arrow; the statement concerns how its components mix.

This correspondence is not accidental. A three-vector and a spin-one state
carry equivalent three-dimensional representations of the rotation group.
The words emph{vector particle} summarize that transformation property. They
do not claim that a quantum polarization state is a tiny rigid arrow hidden
inside the particle.

\section{Discrete and Continuous Symmetries}

Before turning to spin one-half, it helps to distinguish two broad kinds of
groups. A reflection gives a simple discrete example. There is the identity
and there is the reflection \(P\); one either flips or does not flip. By
contrast, a phase
\begin{equation}
e^{i\phi}
\end{equation}
depends continuously on \(\phi\), and a spatial rotation depends continuously
on its axis and angle. Such continuous groups are called \emph{Lie groups},
after Sophus Lie; the name is pronounced \emph{Lee}. Continuous symmetries will
be the main concern, although discrete symmetries remain important.

For the two-element reflection group,
\begin{equation}
P^2=1.
\end{equation}
There is no continuously adjustable amount of reflection between \(1\) and
\(P\). A phase group is different: between \(e^{i\phi_1}\) and
\(e^{i\phi_2}\) lie continuously many other phases. Spatial rotations have
the same continuous character, but unlike \(U(1)\) they are non-abelian.

Spin zero gives the simplest representation of the rotation group. A spin-zero
state is unchanged by every rotation, so every group element is represented
by the one-dimensional identity:
\begin{equation}
D^{(0)}(R)=1.
\end{equation}
The group itself still has infinitely many rotations; it is the
representation on the spin-zero state that is trivial.

This distinction between a group and one of its representations is essential.
The rotation group does not collapse to a one-element group whenever it acts
on a scalar. Rather, every one of its elements is represented by the same
number \(1\) on that particular scalar state.

\section{A General Spin-One-Half State}

A spin-one-half system has two basis states. Relative to a chosen \(z\)-axis,
\begin{equation}
\lvert\uparrow\rangle=
\begin{pmatrix}1\\0\end{pmatrix},
\qquad
\lvert\downarrow\rangle=
\begin{pmatrix}0\\1\end{pmatrix}.
\end{equation}
A general state is a two-component spinor,
\begin{equation}
\lvert\psi\rangle=
\begin{pmatrix}\alpha_1\\\alpha_2\end{pmatrix},
\qquad
|\alpha_1|^2+|\alpha_2|^2=1.
\label{eq:l3-spinor}
\end{equation}
The entries are complex. The choice \((\alpha_1,\alpha_2)=(1,0)\) is spin up,
and \((0,1)\) is spin down. A general normalized combination describes a spin
polarized along some other direction.

The words emph{up} and emph{down} have meaning only after an axis has been
chosen. Here that axis is \(z\). A different measuring axis has a different
pair of basis states, and a state that is definite up along one axis generally
appears as a superposition in another basis. Rotating the physical spin or
rotating the reference axes changes the coefficients accordingly.

Two complex amplitudes contain four real parameters. Normalization removes one
and the physically irrelevant overall phase removes another, leaving two real
parameters. That is exactly the number needed to choose a direction on a
sphere. We need not derive the explicit map between
\((\alpha_1,\alpha_2)\) and that direction in order to study rotations.

The overall phase is removed because
\begin{equation}
\begin{pmatrix}\alpha_1\\\alpha_2\end{pmatrix}
\quad\hbox{and}\quad
e^{i\chi}
\begin{pmatrix}\alpha_1\\\alpha_2\end{pmatrix}
\end{equation}
represent the same physical ray. Normalization fixes the total probability,
not a direction in ordinary Euclidean space. After both identifications, the
remaining two parameters may be regarded as the polar and azimuthal angles of
the spin-polarization direction.

When the spin is rotated, its state becomes
\begin{equation}
\alpha_i'=U_{ij}\alpha_j,
\qquad
\begin{pmatrix}\alpha_1'\\\alpha_2'\end{pmatrix}
=
\begin{pmatrix}
U_{11}&U_{12}\\
U_{21}&U_{22}
\end{pmatrix}
\begin{pmatrix}\alpha_1\\\alpha_2\end{pmatrix}.
\label{eq:l3-spinor-transform}
\end{equation}
Thus spatial rotations possess a two-dimensional matrix representation acting
on spinors. There is a global two-valuedness in this representation that we
defer; locally, the matrices have the transformation properties needed below.

The transformed pair \((\alpha_1',\alpha_2')\) must again be a possible
quantum state. Therefore the first question about \(U\) is not its explicit
axis-angle formula but the condition that it preserve the normalization of
every input spinor.

\section{Probability Preservation and Unitary Matrices}

A rotation must map a normalized state to another normalized state:
\begin{equation}
\alpha_i^*\alpha_i=\alpha_i'^*\alpha_i'=1.
\end{equation}
Substituting Equation~\eqref{eq:l3-spinor-transform},
\begin{align}
\alpha_i'^*\alpha_i'
 &=(U_{ij}\alpha_j)^*(U_{ik}\alpha_k)\\
 &=\alpha_j^*U_{ij}^*U_{ik}\alpha_k.
\end{align}
The analogue of the real-vector condition is therefore
\begin{equation}
(U^T)^*U=\boldsymbol 1.
\end{equation}
The transpose followed by complex conjugation is the Hermitian conjugate,
\begin{equation}
U^\dagger=(U^T)^*,
\end{equation}
so the normalization-preserving condition is
\begin{equation}
U^\dagger U=\boldsymbol 1.
\label{eq:l3-unitarity}
\end{equation}

The extra complex conjugation is the only essential change from the real
vector calculation. For a real vector, \(V_iV_i\) is a length. For a spinor,
\(\alpha_i^*\alpha_i\) is a probability. Complex conjugation is required to
make that probability real and nonnegative, so transposition is replaced by
Hermitian conjugation.

Equation~\eqref{eq:l3-unitarity} also gives the inverse immediately:
\begin{equation}
U^{-1}=U^\dagger.
\end{equation}
The inverse rotation therefore preserves the same probability norm. As with
ordinary rotation matrices, products of unitary matrices remain unitary, so
the representation respects group multiplication.

\begin{figure}[htbp]
\centering
\includegraphics[width=0.62\textwidth]{lecture_03_figure_04.png}
\caption{The spinor rotation matrix obeys \(U^\dagger U=\boldsymbol 1\).}
\lectureframe{00:27:41}
\end{figure}

Matrices satisfying Equation~\eqref{eq:l3-unitarity} are unitary. A general
\(2\times2\) complex matrix begins with eight real parameters. Unitarity is a
\(2\times2\) matrix equation that supplies four independent real constraints,
leaving four parameters. A spatial rotation has only three. One more
restriction is needed to isolate the spinor representation of rotations from
all of \(U(2)\).

The four constraints can also be seen by viewing the two columns of \(U\) as
two complex vectors. Each column must have unit norm, providing two real
conditions, and the columns must be orthogonal, providing one complex or two
real conditions. Altogether that is four. The remaining fourth parameter in
\(U(2)\) will turn out to be an overall determinant phase rather than a new
spatial rotation.

\section{Determinant One and \texorpdfstring{\(SU(2)\)}{SU(2)}}

The additional condition is
\begin{equation}
\det U=1.
\label{eq:l3-det-one}
\end{equation}
It is compatible with group multiplication because
\begin{equation}
\det(AB)=\det A\,\det B.
\end{equation}
The product of two determinant-one unitary matrices again has determinant one,
and the identity has determinant one. These matrices form a closed subgroup.

The inverse is also in the subgroup. If \(\det U=1\), then
\begin{equation}
\det U^{-1}=\frac{1}{\det U}=1.
\end{equation}
Thus imposing determinant one does not remove the identity, spoil closure, or
lose inverses. It consistently selects a group with one fewer continuous
parameter.

Unitarity alone implies
\begin{equation}
\det(U^\dagger U)=
(\det U)^*\det U
=|\det U|^2=1.
\end{equation}
It does not force \(\det U\) to be \(+1\) or \(-1\). A general unitary
determinant may be any phase:
\begin{equation}
\det U=e^{i\phi},
\qquad
\det U^\dagger=e^{-i\phi}.
\end{equation}
Condition~\eqref{eq:l3-det-one} fixes this otherwise free phase and removes one
real parameter.

\begin{classroomqa}
\audiencequestion{Does \(U^\dagger U=\boldsymbol 1\), without the determinant-one condition, force \(\det U^\dagger\) to equal \(\det U\) or to be only \(+1\) or \(-1\)?\lecturetimestamp{00:31:21}}
\lecturerresponse{No. The two determinants are complex conjugates. They can be \(e^{i\phi}\) and \(e^{-i\phi}\), whose product is one. The special-unitary restriction then fixes both to one.}
\end{classroomqa}

The unitary \(n\times n\) matrices form \(U(n)\). For \(n=1\), a unitary
matrix is simply the phase \(e^{i\phi}\), giving \(U(1)\). The determinant-one
subgroup is the \emph{special unitary group},
\begin{equation}
SU(n)=\{U\in U(n):\det U=1\}.
\end{equation}
The two-component spinor is acted on by \(SU(2)\), a three-parameter group.

The notation records both restrictions: \(U\) means unitary, while \(S\)
means special, namely determinant one. These special-unitary groups occur
repeatedly because they preserve a complex inner product while excluding the
independent common phase carried by the determinant.

Locally, \(SU(2)\) has the same rotation algebra and multiplication structure
needed to rotate spin-one-half states. Globally, one must remember the
deferred subtlety: \(SU(2)\) is the double cover of the ordinary rotation group
\(SO(3)\), so two elements of \(SU(2)\) correspond to one spatial rotation.
\editorialnote{The double-cover formulation makes precise the two-valuedness mentioned but not developed in the source.}

\section{Infinitesimal Rotations and the Pauli Matrices}

\begin{classroomqa}
\audiencequestion{Where do the spin matrices, or Pauli matrices, enter this construction?\lecturetimestamp{00:33:16}}
\lecturerresponse{They appear when a group element is taken very close to the identity, corresponding to a rotation through a very small angle. They are the generators of those infinitesimal rotations.}
\end{classroomqa}

Begin with
\begin{equation}
U=\boldsymbol 1+\epsilon M+O(\epsilon^2),
\label{eq:l3-small-u-anti}
\end{equation}
where \(\epsilon\) is real and small. Then
\begin{equation}
U^\dagger=\boldsymbol 1+\epsilon M^\dagger+O(\epsilon^2),
\end{equation}
and unitarity gives
\begin{align}
U^\dagger U
 &=\boldsymbol 1+\epsilon(M^\dagger+M)+O(\epsilon^2).
\end{align}
To first order,
\begin{equation}
M^\dagger=-M.
\end{equation}
Thus the matrix \(M\) in Equation~\eqref{eq:l3-small-u-anti} is
anti-Hermitian.

Only terms linear in \(\epsilon\) are being retained. The product
\(\epsilon^2M^\dagger M\) is discarded because it is of second order. This is
the advantage of working close to the identity: nonlinear group conditions
become linear conditions on their generators.

It is more convenient to extract a factor of \(i\):
\begin{equation}
U=\boldsymbol 1+i\epsilon H+O(\epsilon^2).
\label{eq:l3-small-u-herm}
\end{equation}
Complex conjugation changes \(i\) to \(-i\), so
\begin{equation}
U^\dagger=\boldsymbol 1-i\epsilon H^\dagger+O(\epsilon^2).
\end{equation}
Now
\begin{equation}
U^\dagger U
=\boldsymbol 1+i\epsilon(H-H^\dagger)+O(\epsilon^2),
\end{equation}
and first-order unitarity requires
\begin{equation}
H^\dagger=H.
\end{equation}
With the factor of \(i\) exposed, the generator is Hermitian.

Multiplication by \(i\) has not changed the allowed infinitesimal matrices; it
has only changed the convention used to name the generator. Physics usually
chooses Hermitian generators because their eigenvalues are real and because
the associated observables, such as angular momentum, are Hermitian.

The determinant-one condition supplies tracelessness\editorialnote{The spoken determinant expansion briefly garbles the subscripts. The corrected linear term is \(A_{11}+A_{22}=\operatorname{tr}A\); off-diagonal products begin at quadratic order.}. For a small \(2\times2\)
matrix \(A\),
\begin{equation}
\det(\boldsymbol 1+A)
=
\det
\begin{pmatrix}
1+A_{11}&A_{12}\\
A_{21}&1+A_{22}
\end{pmatrix}.
\end{equation}
Expanding,
\begin{align}
\det(\boldsymbol 1+A)
 &=(1+A_{11})(1+A_{22})-A_{12}A_{21}\\
 &=1+A_{11}+A_{22}+O(A^2)\\
 &=1+\operatorname{tr}A+O(A^2).
\end{align}
Applying this to \(A=i\epsilon H\),
\begin{equation}
\det U=1+i\epsilon\,\operatorname{tr}H+O(\epsilon^2).
\end{equation}
Therefore \(\det U=1\) requires
\begin{equation}
\operatorname{tr}H=0.
\end{equation}
The trace condition removes precisely the generator proportional to the
identity matrix. A general Hermitian \(2\times2\) matrix has four real
parameters: two real diagonal entries and one complex off-diagonal entry. Its
trace fixes one real combination, leaving three. This matches the three axes
about which an infinitesimal spatial rotation can be made.

The infinitesimal generators of \(SU(2)\) are thus traceless Hermitian
\(2\times2\) matrices. The three Pauli matrices form a basis:
\begin{equation}
\sigma_x=
\begin{pmatrix}0&1\\1&0\end{pmatrix},
\qquad
\sigma_y=
\begin{pmatrix}0&-i\\i&0\end{pmatrix},
\qquad
\sigma_z=
\begin{pmatrix}1&0\\0&-1\end{pmatrix}.
\end{equation}
Every generator can be written
\begin{equation}
H=a_x\sigma_x+a_y\sigma_y+a_z\sigma_z.
\end{equation}
A small rotation about the \(x\)-axis selects \(\sigma_x\); one about the
\(y\)-axis selects \(\sigma_y\); and a rotation about any other axis selects
the corresponding linear combination. In units with \(\hbar=1\), the spin
operators are
\begin{equation}
S_i=\frac{\sigma_i}{2}.
\end{equation}

For example, an axis halfway between \(x\) and \(y\) uses a generator
proportional to
\begin{equation}
\frac{\sigma_x+\sigma_y}{\sqrt2}.
\end{equation}
The direction of the coefficient vector
\((a_x,a_y,a_z)\) selects the rotation axis, while its magnitude combines with
\(\epsilon\) to set the small angle. This is the connection between the spin
matrices encountered in angular-momentum commutators and finite rotation
operators. In particular,
\begin{equation}
[S_x,S_y]=iS_z
\end{equation}
and its cyclic partners encode the noncommutativity already seen geometrically.
The standard finite rotation follows by exponentiating the infinitesimal
generator\editorialnote{The exponential finite-rotation formula is the standard completion of the infinitesimal construction; the source derives only the generator conditions and their Pauli-matrix basis.},
\begin{equation}
U(\hat{\boldsymbol n},\theta)
=\exp\left(-\frac{i\theta}{2}\,
\hat{\boldsymbol n}\cdot\boldsymbol\sigma\right).
\end{equation}
\section{A Three-State Color Space}

Combining representations can wait. We now leave spatial spin and return to
quark color. A quark still has spin one-half, but its color is an independent
three-state degree of freedom:
\begin{equation}
\lvert r\rangle,\qquad
\lvert g\rangle,\qquad
\lvert b\rangle.
\end{equation}
A general color state is
\begin{equation}
\lvert q\rangle_{\mathrm{color}}
=
\alpha_r\lvert r\rangle
+\alpha_g\lvert g\rangle
+\alpha_b\lvert b\rangle,
\qquad
|\alpha_r|^2+|\alpha_g|^2+|\alpha_b|^2=1.
\end{equation}
Equivalently,
\begin{equation}
q=
\begin{pmatrix}
\alpha_r\\
\alpha_g\\
\alpha_b
\end{pmatrix}.
\end{equation}
The squared magnitudes are the probabilities for red, green, and blue in a
chosen color basis. The same triplet notation can assemble three quark field
operators,
\begin{equation}
\psi=
\begin{pmatrix}
\psi_r\\
\psi_g\\
\psi_b
\end{pmatrix}.
\end{equation}

Spin and color must be kept conceptually separate. The spin label tells us how
the state transforms under rotations of physical space. The color label tells
us how it transforms under an internal symmetry that does not rotate the
laboratory. A quark state carries both labels, even though the present
discussion suppresses spin in order to concentrate on color.

Writing three field operators in one column is more than bookkeeping. It makes
it possible to state the symmetry in one matrix equation. A transformation
may turn a red field into a linear combination of red, green, and blue fields,
just as a spatial rotation turns one vector component into a linear
combination of the other components.

One can first imagine discrete transformations that exchange two colors. The
larger symmetry contains continuous mixtures of all three. A simple
interchange matrix may need an accompanying phase to make its determinant
equal to one; the exchange is then embedded in the special-unitary group.

For two basis states the bare exchange is represented by
\begin{equation}
X=
\begin{pmatrix}0&1\\1&0\end{pmatrix},
\qquad
\det X=-1.
\end{equation}
Multiplying the whole matrix by \(i\) gives
\begin{equation}
\det(iX)=i^2\det X=1.
\end{equation}
The discrete exchange can therefore appear as a particular element of the
larger continuous group once the phase is chosen consistently.

\begin{classroomqa}
\audiencequestion{Does a matrix that merely interchanges two basis states have the wrong determinant for \(SU(2)\)?\lecturetimestamp{00:43:46}}
\lecturerresponse{The bare interchange has determinant \(-1\). An appropriate factor of \(i\) supplies a determinant-one representative, so the exchange can be included within the continuous special-unitary transformations.}
\end{classroomqa}

\section{\texorpdfstring{\(SU(3)\)}{SU(3)} and the Eight Color Generators}

The continuous color transformation is
\begin{equation}
q'=Uq,
\qquad
U^\dagger U=\boldsymbol 1,
\qquad
\det U=1,
\qquad
U\in SU(3).
\label{eq:l3-su3-action}
\end{equation}
The mathematics is the same special-unitary construction used for spinors,
but now \(U\) is \(3\times3\). This is not a rotation of physical space. It is
a rotation in the internal color space.

Unitarity preserves the color-state normalization:
\begin{equation}
q^\dagger q
=|\alpha_r|^2+|\alpha_g|^2+|\alpha_b|^2.
\end{equation}
Determinant one removes the common \(U(1)\) phase direction from the continuous
mixing. The resulting matrices can perform simple exchanges, relative phase
changes, and genuinely continuous superpositions while remaining in one
closed group.

The parameter count is immediate. A \(3\times3\) complex matrix has nine
complex entries, or eighteen real parameters. The matrix equation
\(U^\dagger U=\boldsymbol 1\) imposes nine real conditions. Determinant one
removes one more:
\begin{equation}
\dim SU(3)=18-9-1=8.
\end{equation}
Thus \(SU(3)\) has eight independent infinitesimal generators.

The nine real unitarity conditions may again be understood in terms of the
columns of \(U\). Three complex columns must each have unit norm, and every
pair must be orthogonal. These conditions remove nine of the original eighteen
real parameters. The determinant phase removes the tenth, leaving eight.
This is the three-state analogue of the \(8-4-1=3\) count for \(SU(2)\).

The appearance of eight is the same fact that ultimately produces eight gluon
species. Saying that there are not nine gluons because the determinant
restriction removed one direction is only a mnemonic. More precisely, the
nine color-anticolor matrices decompose into one color-singlet direction,
proportional to the identity, and eight traceless directions. The eight
traceless generators form the adjoint color structure carried by the QCD
gluons.\editorialnote{The singlet-plus-octet formulation is the standard precise version of the source's intentionally heuristic determinant remark; its gauge-field construction is deferred.}

At this stage the relation to eight gluons is a signpost, not yet a complete
derivation. To obtain gauge particles one must promote the color symmetry to a
local one and determine the compensating fields. What the parameter count has
already shown is that there are eight independent directions in which an
\(SU(3)\) color transformation can begin to move away from the identity.

\section{Color Symmetry Constrains QCD}

The basic symmetry statement is that the physics is invariant under the color
transformation in Equation~\eqref{eq:l3-su3-action}. This is the central
internal symmetry of quantum chromodynamics. Mixing the color components is
analogous in form to mixing spin-up and spin-down amplitudes, and those spinor
transformations are in turn related to rotations of spatial directions, but
the three spaces must not be confused.

Calling color a symmetry means that changing the color basis does not change
physical predictions. Red, green, and blue are convenient names for three
basis vectors, not three observable pigments. A simultaneous \(SU(3)\)
rotation of the color components must leave the form of the laws unchanged.
This invariance is not an optional aesthetic feature; it is the organizing
principle of QCD.

Symmetry restricts the Lagrangian. In the earlier \(U(1)\) example, invariance
under a common phase transformation constrained allowed terms to contain
balanced factors of a field and its conjugate. In QCD, every term must likewise
be invariant under the \(SU(3)\) color transformation. That requirement has
consequences for interactions, conserved quantities, and the form of the
gluon field.

Suppose a proposed term carries an exposed color index. Under
Equation~\eqref{eq:l3-su3-action} that index rotates, so the term cannot by
itself be invariant. Allowed terms must contract color indices into invariant
combinations or use gauge fields whose transformations compensate the change.
This is the practical force of the symmetry principle: it rules out arbitrary
couplings and organizes the interactions that remain.

The full construction of QCD requires one further step: the color
transformation must be allowed to vary from point to point, and the gauge
field must compensate that local change.\editorialnote{Localizing the color symmetry is the standard gauge-theory completion of the global \(SU(3)\) discussion; it is not developed in this source lecture.} That development belongs to the next
stage. For now the essential structure is in place: a color triplet acted on
by a special-unitary three-state symmetry with eight generators.

The calculation has gone as far into group theory as is needed for the moment.
We began with rotations that can be drawn in ordinary space, passed to the
matrices that act on vectors and spinors, and then reused the same logic in an
internal color space. The next step can return to particle physics with the
mathematical framework already available, rather than asking the formalism to
carry the entire physical story at once.
